% ===========================================================================
\section{Löwenheim--Skolem 定理与斯科伦悖论}\label{sec:ls}

完备性定理与紧致性定理直接导致两个表面上与直觉相悖的结论：
\emph{自然数理论有不可数模型，实数理论有可数模型}。二者合起来表明，
一阶公理无法控制其模型的基数。

\begin{theorem}[上行：Peano 公理有不可数模型]\label{thm:upward}
存在 $\M\models\PA$ 使 $|\M|>|\N|$。
\end{theorem}
\begin{proof}
任取不可数集合 $S$。扩充语言：对每个 $s\in S$ 引入常元 $c_s$；
扩充公理：
\[
  c_s\neq c_{s'}\qquad(s\neq s',\ s,s'\in S).
\]
记所得句子集为 $\PA^{+}$。$\PA^{+}$ 的任意有限子集只涉及有限多个 $c_s$，
在 $\N$ 中将这有限个常元解释为互不相同的自然数，即得该有限子集的模型。
由紧致性定理，$\PA^{+}$ 有模型 $\M$。（等价地：矛盾推导只涉及有限条公理，
有限部分一致故 $\PA^{+}$ 一致，再应用定理~\ref{thm:completeness}。）
诸 $c_s^{\M}$ 两两不同，故 $|\M|\ge|S|>\aleph_0$。
\end{proof}

此类 $\M$ 中除标准元素 $0,S(0),S(S(0)),\ldots$ 外还含非标准元素——
例如任何 $c_s^{\M}$ 都不能由闭项 $S(\cdots S(0)\cdots)$ 表出。

\begin{theorem}[下行：可数语言的一致理论有可数模型]\label{thm:downward}
设语言至多可数，$T$ 一致，则 $T$ 有至多可数的模型。
\end{theorem}
\begin{proof}
由注~\ref{rem:countable}，定理~\ref{thm:completeness} 所构造的亨金模型
在语言至多可数时至多可数。应用于实数：实数的一阶理论 $T$ 以
$(\R;0,1,+,\cdot,<)$ 为模型，故一致，从而存在至多可数的 $\M\models T$。
换言之，\emph{实数有可数解释}。
\end{proof}

\begin{remark}[斯科伦悖论]
将下行定理用于实数理论（集合论情形类似），表面上出现矛盾：$\M$ 可数，
而 Cantor 对角线论证可在系统内形式化，因此“$\R$ 不可数”这句一阶句子
在 $\M$ 中为真。矛盾的错觉源于混淆两个层面：
\begin{itemize}[leftmargin=1.3em,itemsep=1pt]
  \item \keyword{系统内}：“不存在从 $\N$ 到 $\R$ 的满射”是 $\M$ 所满足的一阶句子，
        其含义仅为“$\M$ 中不存在这样的函数对象”；
  \item \keyword{系统外}：从外部考察，$\M$ 的论域确实可数，$\N$ 与 $\M$ 所解释的
        “实数”之间确实存在双射——但该双射是系统外的对象，语言中没有符号
        可以指称它，系统内也就无法陈述“此双射存在”。
\end{itemize}
因此，可数与不可数并非一阶语言能够刻画的绝对性质，而只是相对于模型而言的性质。
一阶公理既无法排除不可数的“自然数”，也无法排除可数的“实数”。
\end{remark}

\begin{remark}
若要唯一刻画 $\N$（在同构意义下），需要超出一阶逻辑的手段，例如二阶语义
（对论域的所有子集作量化）或 $\omega$-规则。另一方面，非标准模型并非只是障碍：
非标准分析正是以 $\R$ 的含无穷小元素的扩充模型为工具建立的。
\end{remark}
